\reading{LTR-1}{Liquidity Risk}{STRIKE FIRST}{6}

\begin{lrkeybox}[Learning objectives — official 2026 spine wording]
\begin{enumerate}[label=\textbf{\alph*.}]
\item Explain and calculate liquidity trading risk via cost of liquidation and liquidity-adjusted VaR (LVaR).
\item Identify examples of liquidity funding risk, funding sources, and lessons learned from real cases: Northern Rock, Ashanti Goldfields, and Metallgesellschaft.
\item Evaluate Basel III liquidity risk ratios and BIS principles for sound liquidity risk management.
\item Explain liquidity black holes and identify the causes of positive feedback trading.
\end{enumerate}
{\footnotesize\color{FRMGrey}Source: John C. Hull, \textit{Risk Management and Financial Institutions}, 5th Edition, Chapter 24.}
\end{lrkeybox}

% =====================================================================
\lo{LTR-1 a}{Explain and calculate liquidity trading risk via cost of liquidation and liquidity-adjusted VaR (LVaR).}

\begin{lrdefbox}[Liquidity trading risk]
The risk of moving the price of an asset adversely in the act of buying or selling it.
\end{lrdefbox}

Transaction liquidity risk is low if assets can be liquidated, or a position covered, \term{quickly, cheaply, and without moving the price too much}. A market is said to be \term{liquid} if market participants can put on or unwind positions quickly, without excessive transactions costs and without excessive price deterioration.

\subsection{Bid-offer spread as a measure of liquidity}

\begin{itemize}
\item Higher \term{OFFER} price: the price at which a trader \emph{sells} to you.
\item Lower \term{BID} price: the price at which a trader \emph{buys} from you.
\item The difference between the bid and ask prices is a cost of liquidity.
\end{itemize}

\begin{lrfmlbox}[Dollar and proportional bid-offer spread]
\[ p = \text{Offer price} - \text{Bid price} \qquad\qquad
   s = \frac{\text{Offer price} - \text{Bid price}}{\text{Mid-market price}} \]
\vk{$p$ = dollar bid-offer spread; $s$ = proportional bid-offer spread, a pure number; the mid-market price is halfway between the bid and the offer and can be regarded as the fair price.}
\end{lrfmlbox}

In liquidating a position an institution incurs a cost equal to $s\alpha/2$, where $\alpha$ is the dollar mid-market value of the position. The halving is there because a trade is done on one side of the mid only: a buy at the offer, a sell at the bid.

\subsection{Cost of liquidation}

\begin{lrfmlbox}[Cost of liquidation — normal and stressed markets]
\[ \text{Cost of liquidation (normal market)} \;=\; \sum_{i=1}^{n} \frac{s_i\,\alpha_i}{2} \]
\[ \text{Cost of liquidation (stressed market)} \;=\; \sum_{i=1}^{n} \frac{(\mu_i + \lambda\sigma_i)\,\alpha_i}{2} \]
\vk{$s_i$ = estimated proportional bid-offer spread in normal conditions for instrument $i$; $\alpha_i$ = dollar value of the position in instrument $i$; $n$ = number of positions; $\mu_i$ and $\sigma_i$ = mean and standard deviation of the \emph{proportional} bid-offer spread for instrument $i$; $\lambda$ = the confidence parameter for the spread. For a worst-case spread exceeded only 1\% of the time, with normally distributed spreads, $\lambda = 2.326$.}
\end{lrfmlbox}

\begin{lrtrapbox}[Three polarity and input traps in this formula]
\begin{itemize}
\item \textbf{$\lambda$ is applied to the spread distribution, not the price distribution.} It is the confidence level for how wide the spread gets, not for how far the price moves. That is what separates this from VaR.
\item \textbf{$\mu_i$ and $\sigma_i$ are proportional, not dollar.} The source notes state $\mu$ and $\sigma$ in dollars first and then convert. Using the dollar figures directly against $\alpha_i$ inflates the answer by the price level.
\item \textbf{Equation (1.2) assumes spreads in all instruments are perfectly correlated.} This looks overly conservative and is not: when liquidity tightens, spreads widen together. Diversification reduces market risk; it does not necessarily reduce liquidity trading risk.
\end{itemize}
\end{lrtrapbox}

\gapbadge
\gapnote{The source notes give both formulas but pose no worked example. The following two examples are written from the GARP curriculum (Hull, Chapter 24, Examples 1.1 and 1.2).}

\begin{lrexbox}[A two-instrument book, normal market]
A financial institution has bought 10 million shares of one company and 50 million ounces of a commodity. The shares are bid \$89.5, offer \$90.5. The commodity is bid \$15.00, offer \$15.10. What is the cost of liquidation in a normal market?

\tcblower
\textbf{Step 1 — mid-market prices.}
\[ \text{Shares: } \tfrac{89.5+90.5}{2} = 90.00 \qquad \text{Commodity: } \tfrac{15.00+15.10}{2} = 15.05 \]

\textbf{Step 2 — dollar value of each position, $\alpha_i$.}
\[ \alpha_{\text{shares}} = 90 \times 10 = \$900\text{ million} \qquad
   \alpha_{\text{comm}} = 15.05 \times 50 = \$752.50\text{ million} \]

\textbf{Step 3 — proportional spreads, $s_i$.}
\[ s_{\text{shares}} = \frac{1.00}{90} = 0.01111 \qquad
   s_{\text{comm}} = \frac{0.10}{15.05} = 0.006645 \]

\textbf{Step 4 — apply equation (1.1).}
\[ \frac{900 \times 0.01111}{2} + \frac{752.5 \times 0.006645}{2} = 5.00 + 2.50 = \mathbf{\$7.5\text{ million}} \]

\textbf{What this means.} \$7.5 million is what it costs to walk this book out of the door on a quiet day — half the spread, paid on every dollar of position. It is a cost, not a loss from price movement; the mid-market price has not moved at all in this calculation.
\end{lrexbox}

\begin{lrexbox}[The same book, stressed market]
For the same positions, the mean and standard deviation of the \emph{dollar} bid-offer spread for the shares are \$1.0 and \$2.0. For the commodity both are \$0.1. Spreads are normally distributed. Find the cost of liquidation that will not be exceeded with 99\% confidence.

\tcblower
\textbf{Step 1 — convert the spread moments to proportional terms} by dividing by the mid-market price.
\[ \mu_{\text{shares}} = \tfrac{1.0}{90} = 0.01111 \qquad \sigma_{\text{shares}} = \tfrac{2.0}{90} = 0.02222 \]
\[ \mu_{\text{comm}} = \sigma_{\text{comm}} = \tfrac{0.1}{15.05} = 0.006645 \]

\textbf{Step 2 — set $\lambda$.} A 1\% worst case under normality gives $\lambda = 2.326$.

\textbf{Step 3 — the stressed proportional spread for each instrument.}
\[ \text{Shares: } 0.01111 + 2.326 \times 0.02222 = 0.06279 \]
\[ \text{Commodity: } 0.006645 + 2.326 \times 0.006645 = 0.02210 \]

\textbf{Step 4 — apply equation (1.2).}
\[ 0.5 \times 900 \times 0.06279 \;+\; 0.5 \times 752.5 \times 0.02210 \;=\; 28.26 + 8.32 = \mathbf{\$36.58\text{ million}} \]

\textbf{What this means.} Almost \textbf{five times} the normal-market cost, from the same positions, with no change in price whatsoever. The entire difference is the spread widening. That multiple is the point of the objective: liquidity trading risk is not a rounding error on a stressed day.
\end{lrexbox}

\begin{center}
\begin{tikzpicture}
\begin{axis}[
  width=11.5cm, height=6.6cm,
  ybar stacked, bar width=34pt, enlarge x limits=0.42,
  symbolic x coords={Normal market, Stressed market},
  xtick=data,
  ymin=0, ymax=44,
  ylabel={Cost of liquidation (\$ million)},
  ylabel style={font=\small\sffamily},
  tick label style={font=\small\sffamily},
  ymajorgrids, grid style={FRMRule, dashed},
  legend style={font=\scriptsize\sffamily, draw=FRMRule, fill=white,
    at={(0.5,-0.18)}, anchor=north, legend columns=-1, column sep=10pt},
  area legend,
]
\addplot[fill=PaleNavy, draw=FRMNavy] coordinates {(Normal market,5.00) (Stressed market,28.26)};
\addlegendentry{Shares position}
\addplot[fill=PaleGold, draw=FRMGold] coordinates {(Normal market,2.50) (Stressed market,8.32)};
\addlegendentry{Commodity position}
\node[font=\small\sffamily\bfseries, fill=white, inner sep=1.5pt] at (axis cs:Normal market,10.5) {\$7.5m};
\node[font=\small\sffamily\bfseries, fill=white, inner sep=1.5pt] at (axis cs:Stressed market,40.0) {\$36.58m};
\end{axis}
\end{tikzpicture}
\figcap{The same book, the same prices, two spread assumptions. Stressing the spread alone multiplies the liquidation cost by roughly 4.9$\times$. Note that the shares dominate the increase: their spread volatility ($\sigma = 0.02222$) is twice their mean spread, while the commodity's mean and standard deviation are equal, so the shares absorb far more of the $\lambda\sigma$ term.}
\end{center}

\subsection{Liquidity-adjusted VaR}

VaR estimates a worst-case change in the mark-to-market valuation of the trading book. The cost-of-liquid\-ation measures estimate the cost of liquidating a book \emph{if market prices do not change}. They deal with different risks, so some researchers combine them.

\begin{lrfmlbox}[Liquidity-adjusted VaR]
\[ \text{LVaR (normal market)} \;=\; \text{VaR} \;+\; \sum_{i=1}^{n}\frac{s_i\,\alpha_i}{2} \]
\[ \text{LVaR (stressed market)} \;=\; \text{VaR} \;+\; \sum_{i=1}^{n}\frac{(\mu_i + \lambda\sigma_i)\,\alpha_i}{2} \]
\vk{Notation as above. LVaR is regular VaR \emph{plus} the cost of unwinding — the two are added, never averaged or maximised.}
\end{lrfmlbox}

\begin{lrtrapbox}[The addition is unconditional]
LVaR adds the full liquidation cost to the full VaR. No correlation adjustment, no diversification benefit between the two terms. A distractor that takes $\sqrt{\text{VaR}^2 + \text{cost}^2}$ is built from exactly the instinct a Part 1 candidate brings with them.
\end{lrtrapbox}

% =====================================================================
\lo{LTR-1 b}{Identify examples of liquidity funding risk, funding sources, and lessons learned from real cases: Northern Rock, Ashanti Goldfields, and Metallgesellschaft.}

\begin{lrdefbox}[Liquidity funding risk]
The risk relating to a financial institution's ability to meet its cash needs as they arise. The key to managing it is predicting cash needs and ensuring they can be met in adverse scenarios.
\end{lrdefbox}

\subsection{The main sources of liquidity for a financial institution}

\begin{itemize}
\item Holdings of cash and Treasury securities.
\item The ability to liquidate trading book positions.
\item The ability to borrow money at short notice.
\item The ability to offer favourable terms to attract retail and wholesale deposits at short notice.
\item The ability to securitize assets (such as loans) at short notice.
\item Borrowings from the central bank.
\end{itemize}

\subsection{The three cases}

\begin{center}
\small
\begin{tabularx}{\textwidth}{@{}p{3.05cm} X X@{}}
\toprule
\rowcolor{FRMSoft}\textbf{Case} & \textbf{What they did} & \textbf{Why the cash ran out} \\
\midrule
\textbf{Northern Rock} & Relied on short-term funding for long-term loans (mortgages). & When lenders became wary after the 2007 crisis, it could not access cash and faced a bank run. \\
\addlinespace
\textbf{Ashanti Goldfields} & Used gold forwards to hedge against falling gold prices. & Gold prices \emph{rose}. It could not meet the margin calls because the underlying asset (gold) was illiquid and not easily sold quickly. \\
\addlinespace
\textbf{Metallgesellschaft} & Sold long-dated fixed-price oil supply contracts, hedged with short-dated futures rolled forward. & Oil prices \emph{fell}, producing margin calls on the futures it could not meet. The long-dated contracts that offset them produced no cash. \\
\bottomrule
\end{tabularx}
\end{center}
\figcap{The shared mechanism, and the one worth stating aloud: in all three the \emph{economics} were arguably sound and the \emph{cash timing} was not. Ashanti and Metallgesellschaft are a matched pair in which the price moved in opposite directions, and GARP builds distractors by swapping them.}

\begin{lrtrapbox}[Polarity: which way did the price move?]
Ashanti was hurt by a \textbf{rising} gold price; Metallgesellschaft by a \textbf{falling} oil price. Both were hedgers facing margin calls, so the direction word is the only discriminator in a four-statement audit. State the direction and the actor out loud before locking any answer on these two.
\end{lrtrapbox}

\begin{lrkeybox}[The lesson GARP draws]
The lesson is \emph{not} that firms should avoid hedging. It is that they should ensure they have access to funding to handle the cash flow mismatches that can arise in extreme circumstances.
\end{lrkeybox}

% =====================================================================
\lo{LTR-1 c}{Evaluate Basel III liquidity risk ratios and BIS principles for sound liquidity risk management.}

Basel III introduced two liquidity risk requirements: the \term{liquidity coverage ratio (LCR)} and the \term{net stable funding ratio (NSFR)}.

\begin{enumerate}
\item \textbf{LCR} ensures banks have enough high-quality liquid assets to survive a 30-day stressed period.
\item \textbf{NSFR} focuses on long-term funding stability, requiring banks to have enough stable funding to cover their assets.
\end{enumerate}

\gapbadge
\gapnote{The source notes name both ratios but print neither formula and neither implementation detail. The following is written from the GARP curriculum (Hull, Chapter 24).}

\begin{lrgapbox}
The objective says \emph{evaluate} the ratios. That requires the ratio definitions, what sits in each numerator and denominator, and the stress scenario that defines the LCR's 30-day window.
\end{lrgapbox}

\begin{lrfmlbox}[The two Basel III liquidity ratios]
\[ \text{LCR} \;=\; \frac{\text{High-quality liquid assets}}{\text{Net cash outflows in a 30-day period}} \;\geq\; 100\% \]
\[ \text{NSFR} \;=\; \frac{\text{Amount of stable funding}}{\text{Required amount of stable funding}} \;\geq\; 100\% \]
\vk{NSFR numerator: each category of funding (capital, wholesale deposits, retail deposits, and so on) multiplied by an \term{available stable funding (ASF)} factor reflecting its stability. NSFR denominator: assets and off-balance-sheet items requiring funding, each multiplied by a \term{required stable funding (RSF)} factor.}
\end{lrfmlbox}

\begin{lrkeybox}[The LCR's 30-day stress scenario, and the implementation dates]
The 30-day period in the LCR is one of \textbf{acute stress} involving: a downgrade of three notches (for example AA+ to A+), a partial loss of deposits, a complete loss of wholesale funding, increased haircuts on secured funding, and drawdowns on lines of credit.

LCR was phased in between 2015 and 2019 — 60\% required in 2015, 100\% in 2019. The NSFR implementation date was 1 January 2018.
\end{lrkeybox}

\begin{lrtrapbox}[ASF versus RSF: which side is which]
\term{ASF} weights sit in the \textbf{numerator} and are applied to \textbf{funding} (liabilities and capital) — the more stable the funding, the higher the factor. \term{RSF} weights sit in the \textbf{denominator} and are applied to \textbf{assets} — the harder the asset is to liquidate, the more stable funding it requires. Swapping the two is a textbook 5B.2 sibling swap and it is offered on both papers.
\end{lrtrapbox}

\subsection{BIS principles for sound liquidity risk management}

\begin{itemize}
\item Cover 17 areas for effective liquidity risk management by banks.
\item Emphasise strong governance, risk tolerance, liquidity measurement, diversification of funding sources, and stress testing.
\item Require a contingency funding plan and high-quality liquid assets.
\item Mandate clear public disclosures by banks and comprehensive assessments by supervisors.
\end{itemize}

\begin{lrkeybox}[The first principle, and it is a role question]
A bank is \textbf{responsible for the sound management of its own liquidity risk}, and should establish a robust liquidity risk management framework that ensures it maintains sufficient liquidity, including a cushion of unencumbered, high-quality liquid assets. The supervisor assesses; the bank owns.
\end{lrkeybox}

% =====================================================================
\lo{LTR-1 d}{Explain liquidity black holes and identify the causes of positive feedback trading.}

\begin{lrdefbox}[Liquidity black hole]
A situation where liquidity has dried up in a particular market because everyone wants to sell and no one wants to buy, or vice versa. Sometimes called a \term{crowded exit}.
\end{lrdefbox}

\begin{center}
\small
\begin{tabularx}{\textwidth}{@{}p{4.3cm} X X@{}}
\toprule
\rowcolor{FRMSoft} & \textbf{Negative feedback traders} & \textbf{Positive feedback traders} \\
\midrule
\textbf{Behaviour} & Buy when prices fall, sell when prices rise & Sell when prices fall, buy when prices rise \\
\textbf{Style label} & Value investing & Momentum investing \\
\textbf{Effect on the market} & \textbf{Stabilising} & \textbf{Destabilising} \\
\textbf{When they dominate} & The market is liquid & Prices are liable to be unstable; the market may become one-sided and illiquid \\
\bottomrule
\end{tabularx}
\end{center}

\begin{center}
\begin{tikzpicture}
\begin{axis}[
  width=12.2cm, height=6.2cm,
  xlabel={Time}, ylabel={Price},
  xlabel style={font=\small\sffamily}, ylabel style={font=\small\sffamily},
  tick label style={font=\scriptsize\sffamily},
  xmin=0, xmax=10, ymin=55, ymax=118,
  ymajorgrids, grid style={FRMRule, dashed},
  legend style={font=\scriptsize\sffamily, draw=FRMRule, fill=white,
    at={(0.5,-0.22)}, anchor=north, legend columns=1, row sep=1pt},
]
\addplot[FRMGreen, thick, sharp plot] coordinates
  {(0,100)(1,96)(2,99)(3,94)(4,98)(5,95)(6,99)(7,96)(8,100)(9,97)(10,99)};
\addlegendentry{Negative feedback dominates: buying into falls damps the move}
\addplot[FRMRed, thick, sharp plot] coordinates
  {(0,100)(1,96)(2,92)(3,86)(4,82)(5,76)(6,71)(7,67)(8,63)(9,60)(10,58)};
\addlegendentry{Positive feedback dominates: selling into falls feeds the move}
\node[font=\scriptsize\sffamily, fill=white, inner sep=1.5pt, anchor=west] at (axis cs:5.4,104) {market stays two-sided};
\node[font=\scriptsize\sffamily, fill=white, inner sep=1.5pt, anchor=west] at (axis cs:4.1,66.5) {crowded exit};
\draw[-{Stealth[length=4pt]}, FRMRed] (axis cs:6.1,65.8) -- (axis cs:8.8,61.0);
\end{axis}
\end{tikzpicture}
\figcap{The same initial shock, two trader populations. The mechanism is not that positive feedback traders are wrong about value — it is that their reaction function has the same sign as the price move, so each round of selling manufactures the next. That is what turns a price fall into a black hole.}
\end{center}

\begin{lrkeybox}[Causes of positive feedback trading]
\begin{itemize}
\item Trend trading
\item Stop-loss rules
\item Dynamic hedging
\end{itemize}
All three share one feature worth naming: each \emph{mechanically} requires selling as the price falls. None of them is a judgement about value.
\end{lrkeybox}

\begin{lrtrapbox}[Consolidated trap summary — LTR-1]
\begin{itemize}
\item \textbf{Polarity.} Negative feedback = buy on falls = stabilising. Positive = sell on falls = destabilising. Reversing this pair is the single most likely distractor in this reading.
\item \textbf{Polarity.} Ashanti = gold price \emph{rose}; Metallgesellschaft = oil price \emph{fell}.
\item \textbf{Sibling.} ASF (funding, numerator) versus RSF (assets, denominator).
\item \textbf{Input twin.} $\mu$ and $\sigma$ in equation (1.2) are \emph{proportional} spreads, not dollar spreads.
\item \textbf{Sign.} LVaR = VaR \emph{plus} liquidation cost. Straight addition.
\item \textbf{Scope.} $\lambda$ is the confidence level on the \emph{spread} distribution, not on the price distribution.
\item \textbf{Intermediate result.} In a two-instrument cost-of-liquidation question, each instrument's own contribution (\$5.00m, \$2.50m, \$28.26m, \$8.32m here) is a natural offered distractor. Read whether the question asks for the book or for one leg.
\end{itemize}
\end{lrtrapbox}
